\section{Unit tests of the structural signals}
\label{app:tests}
All properties used in \cref{sec:method} are checked by automated tests that ship with the code
(35 tests, all passing; the list covers the properties used in this paper).
\begin{itemize}
\item \textbf{Exact Euler characteristic.} On toy shapes (disk, annulus, two disks, a lens-shaped
  leaf, and the lens with two holes; $\chi=1,0,2,1,-1$) the cubical $\chi=V-E+F$ of the binary set
  equals the value computed independently with scikit-image (4-connected foreground).
\item \textbf{Telescoping.} For every direction, $\sum_k\dchi(\theta,k)$, $\sum_k\Delta A(\theta,k)$ and
  $\sum_k\Delta P(\theta,k)$ equal the total valuation of the shape, to numerical precision.
\item \textbf{Partition of unity.} Slab memberships~\eqref{eq:slab} sum to one for every height.
\item \textbf{Sharp limit.} With binary weights and large $\kappa$, the relaxed increments equal those
  of a hard sweep that assigns each cell to the slab containing its height.
\item \textbf{Closed form.} For an $n\times n$ block of pixels, the totals are
  $(A,P,\chi)=((n-1)^2,4(n-1),1)$.
\item \textbf{Consistency.} The Euler column of the Hadwiger increments equals the separately
  implemented Euler increments.
\item \textbf{Rotation.} For unit weights, rotating the input by $90^\circ$ cyclically shifts
  $\dchi$, $\Delta A$ and $\Delta P$ by $K/4$ directions. This follows from
  $\langle \mathcal R x,\theta\rangle=\langle x,\mathcal R^{-1}\theta\rangle$, the equivariance of the
  centroid/RMS frame and the fact that a $90^\circ$ grid rotation maps cells of each dimension
  bijectively onto cells of the same dimension; for rotations by $2\pi/K$ it holds up to grid resampling.
\item \textbf{Mask pooling.} The mask-pooling control is unchanged when patch tokens outside the leaf
  mask are altered.
\item \textbf{Fuzzy subcomplex.} Learned cell weights satisfy $w(\text{coface})\le w(\text{face})$, and
  gradients reach the per-pixel weights.
\item \textbf{Splits.} Under the grouped protocol no exact duplicate or re-shot group crosses a fold
  boundary and every image is tested exactly once; under the random protocol leaks occur.
\end{itemize}

\section{Species and taxonomy}
\begin{table}[h]
\centering
\caption{The 30 species (APG~IV order and family) and the number of images after de-duplication.}
\label{tab:taxonomy}
\scriptsize
\resizebox{\linewidth}{!}{\begin{tabular}{lllr}
\toprule
Order / family & Species & Name & $n$ \\
\midrule
Brassicales / Brassicaceae & \textit{Brassica juncea} & Indian Mustard & 34 \\
Brassicales / Moringaceae & \textit{Moringa oleifera} & Drumstick & 77 \\
Caryophyllales / Amaranthaceae & \textit{Amaranthus viridis} & Arive-Dantu & 122 \\
Caryophyllales / Basellaceae & \textit{Basella alba} & Basale & 103 \\
Fabales / Fabaceae & \textit{Pongamia pinnata} & Indian Beech & 61 \\
Fabales / Fabaceae & \textit{Trigonella foenum-graecum} & Fenugreek & 36 \\
Gentianales / Apocynaceae & \textit{Carissa carandas} & Karanda & 64 \\
Gentianales / Apocynaceae & \textit{Nerium oleander} & Oleander & 62 \\
Gentianales / Apocynaceae & \textit{Tabernaemontana divaricata} & Crape Jasmine & 56 \\
Lamiales / Lamiaceae & \textit{Mentha sp.} & Mint & 97 \\
Lamiales / Lamiaceae & \textit{Ocimum tenuiflorum} & Tulsi & 52 \\
Lamiales / Lamiaceae & \textit{Plectranthus amboinicus} & Mexican Mint & 48 \\
Lamiales / Oleaceae & \textit{Jasminum sp.} & Jasmine & 71 \\
Lamiales / Oleaceae & \textit{Nyctanthes arbor-tristis} & Parijata & 40 \\
Malvales / Malvaceae & \textit{Hibiscus rosa-sinensis} &  & 43 \\
Malvales / Muntingiaceae & \textit{Muntingia calabura} & Jamaica Cherry-Gasagase & 56 \\
Myrtales / Lythraceae & \textit{Punica granatum} & Pomegranate & 79 \\
Myrtales / Myrtaceae & \textit{Psidium guajava} & Guava & 65 \\
Myrtales / Myrtaceae & \textit{Syzygium cumini} & Jamun & 39 \\
Myrtales / Myrtaceae & \textit{Syzygium jambos} & Rose Apple & 56 \\
Piperales / Piperaceae & \textit{Piper betle} & Betel & 48 \\
Rosales / Moraceae & \textit{Artocarpus heterophyllus} & Jackfruit & 56 \\
Rosales / Moraceae & \textit{Ficus auriculata} & Roxburgh fig & 50 \\
Rosales / Moraceae & \textit{Ficus religiosa} & Peepal Tree & 63 \\
Santalales / Santalaceae & \textit{Santalum album} & Sandalwood & 58 \\
Sapindales / Anacardiaceae & \textit{Mangifera indica} & Mango & 62 \\
Sapindales / Meliaceae & \textit{Azadirachta indica} & Neem & 60 \\
Sapindales / Rutaceae & \textit{Citrus limon} & Lemon & 57 \\
Sapindales / Rutaceae & \textit{Murraya koenigii} & Curry & 60 \\
Zingiberales / Zingiberaceae & \textit{Alpinia galanga} & Rasna & 50 \\
\bottomrule
\end{tabular}}
\end{table}
\section{Field-photo test set}
\label{app:field}
\begin{table}[h]
\centering
\caption{DIMPSAR folders used for the field test, the species they were mapped to (checked visually),
and the number of images in which our masker found a leaf. Kannada \emph{Arali} denotes
\textit{Ficus religiosa}, not \textit{Nerium oleander}.}
\label{tab:dimpsar}
\scriptsize
\begin{tabular}{lllr}
\toprule
DIMPSAR folder & Species & Family & Images used \\
\midrule
Jackfruit & \textit{Artocarpus heterophyllus} & Moraceae & 110 \\
Neem & \textit{Azadirachta indica} & Meliaceae & 113 \\
Malabar Spinach & \textit{Basella alba} & Basellaceae & 79 \\
Lemon & \textit{Citrus limon} & Rutaceae & 120 \\
Arali & \textit{Ficus religiosa} & Moraceae & 89 \\
Hibiscus & \textit{Hibiscus rosa-sinensis} & Malvaceae & 113 \\
Jasmine & \textit{Jasminum sp.} & Oleaceae & 48 \\
Mango & \textit{Mangifera indica} & Anacardiaceae & 103 \\
Mint & \textit{Mentha sp.} & Lamiaceae & 122 \\
Drumstick & \textit{Moringa oleifera} & Moringaceae & 26 \\
Gasagase & \textit{Muntingia calabura} & Muntingiaceae & 77 \\
Curry & \textit{Murraya koenigii} & Rutaceae & 159 \\
Parijatha & \textit{Nyctanthes arbor-tristis} & Oleaceae & 61 \\
Tulsi & \textit{Ocimum tenuiflorum} & Lamiaceae & 171 \\
Betel & \textit{Piper betle} & Piperaceae & 112 \\
Doddpathre & \textit{Plectranthus amboinicus} & Lamiaceae & 131 \\
Honge & \textit{Pongamia pinnata} & Fabaceae & 110 \\
Guava & \textit{Psidium guajava} & Myrtaceae & 121 \\
Pomoegranate & \textit{Punica granatum} & Lythraceae & 73 \\
Nerale & \textit{Syzygium cumini} & Myrtaceae & 62 \\
\bottomrule
\end{tabular}
\end{table}

\begin{table}[h]
\centering
\caption{Per-species 20-way accuracy (\%) on raw field photographs, ConvNeXt-T, mean over all runs.}
\label{tab:fieldspecies}
\scriptsize
\resizebox{\linewidth}{!}{\begin{tabular}{lrccccc}
\toprule
Species & $n$ & Plain & Global & Strip & Slab-only & SVA \\
\midrule
\textit{Syzygium cumini} & 62 & 15 & 14 & 8 & 4 & 2 \\
\textit{Nyctanthes arbor-tristis} & 61 & 4 & 4 & 8 & 11 & 4 \\
\textit{Jasminum sp.} & 48 & 2 & 0 & 6 & 6 & 5 \\
\textit{Pongamia pinnata} & 110 & 12 & 8 & 24 & 16 & 17 \\
\textit{Mangifera indica} & 103 & 13 & 16 & 21 & 17 & 20 \\
\textit{Artocarpus heterophyllus} & 110 & 22 & 20 & 34 & 32 & 31 \\
\textit{Moringa oleifera} & 26 & 35 & 33 & 34 & 27 & 36 \\
\textit{Citrus limon} & 120 & 13 & 14 & 25 & 25 & 43 \\
\textit{Mentha sp.} & 122 & 22 & 24 & 54 & 53 & 45 \\
\textit{Murraya koenigii} & 159 & 57 & 53 & 54 & 44 & 49 \\
\textit{Ocimum tenuiflorum} & 171 & 44 & 38 & 50 & 58 & 50 \\
\textit{Hibiscus rosa-sinensis} & 113 & 30 & 33 & 40 & 40 & 50 \\
\textit{Azadirachta indica} & 113 & 16 & 21 & 45 & 33 & 53 \\
\textit{Psidium guajava} & 121 & 41 & 35 & 52 & 53 & 55 \\
\textit{Plectranthus amboinicus} & 131 & 37 & 36 & 66 & 71 & 61 \\
\textit{Basella alba} & 79 & 54 & 56 & 63 & 63 & 62 \\
\textit{Punica granatum} & 73 & 25 & 41 & 67 & 72 & 64 \\
\textit{Muntingia calabura} & 77 & 28 & 28 & 53 & 60 & 65 \\
\textit{Piper betle} & 112 & 77 & 80 & 84 & 82 & 86 \\
\textit{Ficus religiosa} & 89 & 80 & 83 & 91 & 91 & 95 \\
\bottomrule
\end{tabular}}
\end{table}

\section{Computational cost}
\label{app:cost}
\begin{table}[h]
\centering
\caption{Parameters, GFLOPs (one image at training resolution) and median batch-1 latency on a laptop
RTX~3050 GPU and on CPU. Mask pooling and strip pooling add almost no cost; the sweep-based models
roughly double CPU and GPU latency, mostly because of the histogram operations.}
\label{tab:cost}
\scriptsize
\resizebox{\linewidth}{!}{\begin{tabular}{llrrrrr}
\toprule
Backbone & Model & Params (M) & Added (M) & GFLOPs & GPU ms & CPU ms \\
\midrule
ConvNeXt-T & Plain & 27.9 & 0.12 & 27.9 & 16 & 341 \\
 & Mask pooling & 28.1 & 0.23 & 28.0 & 21 & 347 \\
 & Strip pooling & 28.9 & 1.11 & 28.0 & 18 & 318 \\
 & Slab-only & 29.1 & 1.26 & 29.3 & 26 & 661 \\
 & SVA & 29.1 & 1.30 & 29.3 & 30 & 621 \\
 & Global fusion & 29.6 & 1.79 & 28.1 & 28 & 411 \\
\midrule
EfficientNet-B0 & Plain & 4.1 & 0.09 & 3.2 & 23 & 270 \\
 & Strip pooling & 5.0 & 1.03 & 3.2 & 25 & 294 \\
 & Slab-only & 5.1 & 1.12 & 4.3 & 31 & 386 \\
 & SVA & 5.2 & 1.16 & 4.3 & 49 & 431 \\
 & Global fusion & 5.8 & 1.75 & 3.3 & 39 & 348 \\
\midrule
ViT-S/16 & Plain & 22.1 & 0.27 & 29.6 & 16 & 462 \\
 & Strip pooling & 23.1 & 1.27 & 29.7 & 17 & 510 \\
 & Slab-only & 23.2 & 1.42 & 31.2 & 22 & 479 \\
 & SVA & 23.3 & 1.46 & 31.3 & 30 & 635 \\
 & Global fusion & 23.8 & 1.96 & 30.1 & 32 & 569 \\
\midrule
Leaf masker (448$^2$) & -- & 28.0 & -- & 38.2 & 27 & 307 \\
\bottomrule
\end{tabular}}
\end{table}